\documentclass[11pt]{article}
% =====================================================================
%  Status:      research note (exploratory).  NOT a paper claim.
%  Provenance:  derived in the session of 2026-09-28/29; depends on the
%               group conventions fixed in companion note 07.
%  Promotion:   nothing here is promoted.  See §6.
% =====================================================================
\usepackage[utf8]{inputenc}
\usepackage[T1]{fontenc}
\usepackage{amsmath}
\usepackage{amssymb}
\usepackage{amsthm}
\usepackage{booktabs}
\usepackage{hyperref}
\hypersetup{colorlinks=true,linkcolor=blue,urlcolor=cyan}
\usepackage[a4paper,margin=1in]{geometry}

\newtheorem{theorem}{Theorem}[section]
\newtheorem{lemma}[theorem]{Lemma}
\newtheorem{proposition}[theorem]{Proposition}
\newtheorem{corollary}[theorem]{Corollary}
\newtheorem{definition}[theorem]{Definition}
\newtheorem{remark}[theorem]{Remark}
\newtheorem{example}[theorem]{Example}

\tolerance=1000
\emergencystretch=2.5em

\title{Automorphisms of the joint group, dilation, and two obstructions}
\author{Mingli Yuan and DeepSeek Harness Agent}
\date{2026-09-29}

\begin{document}
\maketitle

\begin{abstract}
The joint group $G=\mathbb R_A\times\operatorname{Aff}^+(1,\mathbb R)$ has a
four-parameter family of Lie group automorphisms. Two of the parameters are
geometrically meaningful: the diagonal one $\Phi_\lambda$ realises the
\emph{similarity} of the plane in which the group acts, and the off-diagonal one
$\gamma$ is the only part that cannot be seen in the plane at all. This note
proves the classification, identifies the inner automorphisms, establishes the
compatibility $D_\lambda T_gD_\lambda^{-1}=T_{\Phi_\lambda(g)}$, compares the
outer hierarchy $\Phi_2^n(\Gamma)$ with the inner one and records why only the
former removes both translation directions, and gives two obstructions to
$\gamma\ne0$: the formula is not semialgebraic in the multiplicative coordinates,
and it cannot be implemented by conjugating the plane action.
\end{abstract}

\tableofcontents

\section{Status and provenance}
\label{sec:status}

\textbf{Status.} Research note, exploratory. Theorems~\ref{thm:aut} and the
propositions of \S\S\ref{sec:dilation}--\ref{sec:obstructions} are proved here.
$G$ is not a Paper 0 object and none of this is a paper claim.
Proposition~\ref{prop:conformal} is the smallest item likely to be reusable.

\textbf{Provenance.} Written 2026-09-29. The group law, the centre, the derived
subgroup and the plane action are those of the companion note 07 in this
directory; the contact form and the horizontal metric are Paper 0 \S9.

\section{The automorphism group}
\label{sec:aut}

Recall
\[
G=\{(A,S,B):A,B\in\mathbb R,\ S>0\},\qquad
(A_2,S_2,B_2)(A_1,S_1,B_1)=(A_2+A_1,\ S_2S_1,\ B_2+S_2B_1),
\]
with
\[
\mathfrak g=\operatorname{span}\{Z,H,N\},\quad
[Z,H]=[Z,N]=0,\quad [H,N]=N,
\]
$Z$ the $A$-direction, $H$ the $M=\log S$ direction and $N$ the $B$-direction.
As $G\cong\mathbb R^3$ it is connected and simply connected, so a Lie algebra
automorphism integrates to a unique Lie group automorphism, and conversely.

\begin{theorem}[All Lie automorphisms]
\label{thm:aut}
Every Lie group automorphism of $G$ is of the form
\begin{equation}
\label{eq:F}
F_{\alpha,\gamma,\delta,\kappa}(A,S,B)
=\bigl(\alpha A+\gamma\log S,\ S,\ \delta B+\kappa(S-1)\bigr),
\qquad \alpha\delta\ne0,
\end{equation}
and every map of this form is an automorphism. The parameters compose by
\[
(\alpha_2,\gamma_2,\delta_2,\kappa_2)\circ(\alpha_1,\gamma_1,\delta_1,\kappa_1)
=(\alpha_2\alpha_1,\ \alpha_2\gamma_1+\gamma_2,\ \delta_2\delta_1,\ \delta_2\kappa_1+\kappa_2).
\]
\end{theorem}

\begin{proof}
That \eqref{eq:F} is a homomorphism is the direct computation
\[
F(gh)=\bigl(\alpha(A_2{+}A_1)+\gamma\log(S_2S_1),\ S_2S_1,\
\delta(B_2{+}S_2B_1)+\kappa(S_2S_1{-}1)\bigr)=F(g)F(h),
\]
the third coordinate matching because
$\kappa(S_2S_1-1)=\kappa(S_2-1)+S_2\kappa(S_1-1)$; it is bijective exactly when
$\alpha\delta\ne0$, with $F(g)=e$ only for $g=e$.

Conversely, let $F$ be any automorphism. The centre and the derived subgroup are
characteristic, so $dF$ maps $\mathbb RZ$ to $\mathbb RZ$ and $\mathbb RN$ to
$\mathbb RN$:
\[
dF(Z)=\alpha Z,\qquad dF(N)=\delta N,\qquad dF(H)=\beta H+\gamma Z+\kappa N .
\]
Applying $dF$ to $[H,N]=N$ gives
$\beta\delta N=\delta N$, hence $\beta=1$ because $\delta\ne0$. Since $G$ is
connected and simply connected, $F$ is determined by $dF$, and the differential
just computed is exactly that of \eqref{eq:F}.
\end{proof}

\begin{proposition}[Inner automorphisms]
\label{prop:inner}
Conjugation by $k=(\alpha,r,\beta)$ sends
\[
(A,S,B)\longmapsto\bigl(A,\ S,\ rB-(S-1)\beta\bigr).
\]
Hence the inner automorphisms are exactly the $F_{1,0,r,\kappa}$ with $r>0$ and
$\kappa$ arbitrary.
\end{proposition}

\begin{proof}
$kgk^{-1}=(A,S,rB+(1-S)\beta)$, which in the parametrisation \eqref{eq:F} is
$\alpha=1$, $\gamma=0$, $\delta=r$, $\kappa=-\beta$. Conversely every
$F_{1,0,r,\kappa}$ with $r>0$ is conjugation by $(0,r,-\kappa)$.
\end{proof}

\begin{proposition}[No automorphism cycles the axes]
\label{prop:no-cycle}
There is no automorphism of $G$ interchanging the three coordinate directions
cyclically.
\end{proposition}

\begin{proof}
By Theorem~\ref{thm:aut}, $dF(Z)=\alpha Z$ and $dF(N)=\delta N$ are always
multiples of $Z$ and $N$ respectively, while $dF(H)=H+\gamma Z+\kappa N$ always
has $H$-coefficient exactly $1$. So no generator is sent to a different
generator.
\end{proof}

\section{Dilation is an outer automorphism}
\label{sec:dilation}

Put
\[
\Phi_\lambda(A,S,B)=(\lambda A,\ S,\ \lambda B),\qquad \lambda>0,
\]
which is $F_{\lambda,0,\lambda,0}$ in \eqref{eq:F}. Write also
\[
T_g(x,y)=(x+A,\ Sy+B),\qquad D_\lambda(x,y)=(\lambda x,\ \lambda y).
\]

\begin{proposition}
\label{prop:dilation}
\leavevmode
\begin{enumerate}
\item $\Phi_\lambda\in\operatorname{Aut}(G)$ and
$D_\lambda T_gD_\lambda^{-1}=T_{\Phi_\lambda(g)}$ for every $g$ and every
$\lambda>0$.
\item For $\lambda\ne1$, $\Phi_\lambda$ is \emph{not} inner.
\end{enumerate}
\end{proposition}

\begin{proof}
(1) The first part is Theorem~\ref{thm:aut}. For the conjugation identity,
\[
D_\lambda T_gD_\lambda^{-1}(x,y)
=D_\lambda\Bigl(\frac x\lambda+A,\ \frac{Sy}\lambda+B\Bigr)
=\bigl(x+\lambda A,\ Sy+\lambda B\bigr)=T_{\Phi_\lambda(g)}(x,y).
\]
(2) Inner automorphisms act trivially on the centre, since the centre is central;
but $\Phi_\lambda(A,1,0)=(\lambda A,1,0)\ne(A,1,0)$ for $\lambda\ne1$ and
$A\ne0$.
\end{proof}

So the similarity of the acting plane is implemented by a group automorphism of
the joint quotient, and for $\lambda\ne1$ by one that no conjugation inside $G$
can produce. This is a faithful compatibility between an automorphism of the
group and a geometric transformation of the plane.

\begin{proposition}[Dilation scales the contact form]
\label{prop:conformal}
In the coordinates $(A,M,B)$ with $M=\log S$,
\[
\Phi_\lambda^*\alpha=\lambda\,\alpha,
\qquad
\Phi_\lambda^*\bigl(dA^2+dM^2\bigr)=\lambda^2dA^2+dM^2,
\]
where $\alpha=dB-dA-B\,dM$. Hence $\Phi_\lambda$ rescales the contact form by
the constant $\lambda$ but is an isometry of the horizontal metric only when
$\lambda=1$.
\end{proposition}

\begin{proof}
$\Phi_\lambda$ reads $(A,M,B)\mapsto(\lambda A,M,\lambda B)$, so
$d(\lambda B)-d(\lambda A)-(\lambda B)dM=\lambda(dB-dA-B\,dM)$. The horizontal
metric pulls back to $\lambda^2dA^2+dM^2$.
\end{proof}

\begin{remark}
Consequently $D_\lambda T_gD_\lambda^{-1}=T_{\Phi_\lambda(g)}$ is a statement
about \emph{similarities}, not isometries, of the plane action; and it is
consistent with the companion note's observation that no positive definite
metric makes the whole of $T(G)$ act by isometries.
\end{remark}

\section{Hierarchies: outer versus inner}
\label{sec:hierarchy}

Let
\[
\tau=(1,1,0),\qquad \eta=(0,1,1),\qquad \delta=(0,2,0),
\qquad \Gamma=\{(m,1,n):m,n\in\mathbb Z\}\cong\mathbb Z^2,
\]
so that $\Gamma$ is the subgroup of $G$ acting as the unit square lattice.

\begin{proposition}[Inner hierarchy]
\label{prop:inner-hierarchy}
$\delta\tau\delta^{-1}=\tau$ and $\delta\eta\delta^{-1}=\eta^2$, and
geometrically $T_\delta(x,y)=(x,2y)$ with $T_\delta Q=Q\cup T_\eta Q$ for
$Q=[0,1]^2$. Consequently the conjugates $\Gamma_n=\mathbb Z\times2^n\mathbb Z$
satisfy
\[
\bigcap_{n\ge0}\Gamma_n=\mathbb Z\times\{0\}\ne\{0\}.
\]
\end{proposition}

\begin{proof}
The conjugations are direct. For the geometric statement,
$T_\delta Q=[0,1]\times[0,2]=Q\cup([0,1]\times[1,2])=Q\cup T_\eta Q$, with
shared boundary only. The subgroups $\Gamma_n$ are the stabilisers of the
$n$-fold refined lattices, and the intersection is the set of pairs $(m,0)$.
\end{proof}

So an internal refinement scales only one direction: however deep the vertical
subdivision, the horizontal integer translations survive.

\begin{proposition}[Outer hierarchy removes both]
\label{prop:outer-hierarchy}
$\Phi_2(\tau)=\tau^2$, $\Phi_2(\eta)=\eta^2$, and
\[
\Phi_2^n(\Gamma)=2^n\mathbb Z^2,
\qquad
\bigcap_{n\ge0}\Phi_2^n(\Gamma)=\{0\}.
\]
\end{proposition}

\begin{proof}
$\Phi_2(A,S,B)=(2A,S,2B)$ acts on $\Gamma=\{(m,1,n)\}$ by $(m,n)\mapsto(2m,2n)$.
The intersection of the subgroups $2^n\mathbb Z^2$ is trivial.
\end{proof}

\begin{remark}
The two hierarchies are different mechanisms and must not be conflated. Only the
outer family, which is exactly the family implementing dilations of the plane,
removes a period in \emph{every} direction. This does \emph{not} by itself
produce any tiling statement: whether a given tile geometry \emph{forces} such a
hierarchy is a separate question and is not addressed here.
\end{remark}

\section{The two obstructions to \texorpdfstring{$\gamma\ne0$}{gamma != 0}}
\label{sec:obstructions}

The parameter $\gamma$ in \eqref{eq:F} mixes the centre into the $H$-direction.
It is the only part of $\operatorname{Aut}(G)$ that the geometry of the acting
plane cannot see. Two obstructions make this precise.

\begin{proposition}[$\log$ is not semialgebraic]
\label{prop:log-not-semialgebraic}
The graph $\{(x,y):x>0,\ y=\log x\}$ is not a semialgebraic subset of
$\mathbb R^2$. Consequently no $F_{\alpha,\gamma,\delta,\kappa}$ with
$\gamma\ne0$ has a semialgebraic formula in the multiplicative coordinates
$(A,S,B)$.
\end{proposition}

\begin{proof}
Suppose some nonzero polynomial $P$ vanished on the whole graph. Substituting
$x=e^t$, $y=t$ gives
\[
P(e^t,t)=\sum_{i=0}^{n}p_i(t)e^{it}=0\qquad\text{for all }t\in\mathbb R,
\]
with $p_n\ne0$. Dividing by $e^{nt}$ and letting $t\to+\infty$ forces
$p_n(t)\to0$, which is impossible for a nonzero polynomial. Hence no such $P$
exists, and the graph --- having empty interior --- cannot be described by
finitely many polynomial equalities and inequalities either.
\end{proof}

\begin{proposition}[$\gamma\ne0$ is not a plane transformation]
\label{prop:gamma-not-plane}
No $F_{\alpha,\gamma,\delta,\kappa}$ with $\gamma\ne0$ is implemented by
conjugating the plane action $T(G)$ with a global bijection of the plane.
\end{proposition}

\begin{proof}
For $(x,y)$ in the plane, the stabiliser in $T(G)$ is
\[
\operatorname{Stab}(x,y)=\bigl\{(0,S,(1-S)y):S>0\bigr\},
\]
every element of which has $A$-coordinate $0$. Any conjugation of the action by a
bijection carries stabilisers to stabilisers. But
$F_{\alpha,\gamma,\delta,\kappa}(0,S,(1-S)y)$ has $A$-coordinate
$\gamma\log S$, which is nonzero for $S\ne1$ when $\gamma\ne0$, so the image of
the stabiliser is not the stabiliser of any point. Contradiction.
\end{proof}

\begin{remark}
By contrast, when $\gamma=0$ every automorphism of \eqref{eq:F} \emph{is}
implemented by a plane affine map: conjugating $T(G)$ with
\[
F_{\mathrm{plane}}(x,y)=(\alpha x,\ \delta y-\kappa)
\]
gives
\[
F_{\mathrm{plane}}T_gF_{\mathrm{plane}}^{-1}=T_{F_{\alpha,0,\delta,\kappa}(g)},
\]
as a direct substitution shows.
\end{remark}

\section{Boundaries: what is not claimed}
\label{sec:boundaries}

\begin{enumerate}
\item \textbf{Lie automorphisms only.} Theorem~\ref{thm:aut} classifies the
\emph{smooth} automorphisms of the Lie group. Abstract group automorphisms
without any continuity hypothesis are \emph{not} classified, and no such claim is
made.
\item \textbf{Structure preservation is not automatic.} Theorem~\ref{thm:aut}
gives the automorphisms of the group. It does not say that every one of them
preserves the contact form, the horizontal metric, the assignment, the
projections of the companion note, or any verification certificate.
Proposition~\ref{prop:conformal} is the one case computed here.
\item \textbf{The hierarchy is not a tiling theorem.} Nothing above asserts that
a tile geometry forces an intrinsic hierarchy, or that any particular tile is
aperiodic. The hierarchy statements concern subgroups of $G$ only.
\item \textbf{$\Phi_\lambda$ is not a symmetry of the contact model.} It rescales
$\alpha$ by $\lambda$ and the horizontal metric anisotropically; the maps it
implements on the plane are similarities, not isometries.
\item \textbf{The finite checks are regressions.} The identities above are proved
symbolically; the exact rational family checked alongside them is a regression
only and cannot be quoted as a proof of any universal statement.
\end{enumerate}

\begin{thebibliography}{9}
\bibitem{Paper0}
M.~Yuan and DeepSeek Harness Agent,
\emph{Arithmetic Expression Geometry 0: Arithmetic Expression Spaces},
\path{paper-0/}, commit \texttt{d89ceb8f8090482317b750407f623c8dcfb37ea2}.
\bibitem{Note07}
Companion note 07 in this directory,
\emph{The joint quotient of the ACS charge and the AES affine operator}.
\bibitem{Note08}
Companion note 08 in this directory,
\emph{Isometric realisations of the joint group in the eight Thurston geometries}.
\end{thebibliography}

\end{document}
